\documentclass{zhvt-classic}

% 使用本地源流明体（GenRyuMin2 TW）按文件路径加载，不依赖系统字体名
% 版框（内框／界栏／书口侧栏／鱼尾）改用与版心印章同一印泥棕色：
%   印章图片 yshs.png 实色 = RGB(135,68,52) = HTML 874434
\usepackage{xcolor}
\usepackage{tikz}
% sealbrown 双注册：xcolor 供 TikZ/pgf，L3 供 cls 版框
\definecolor{sealbrown}{HTML}{874434}
\ExplSyntaxOn
\color_set:nnn { sealbrown } { HTML } { 874434 }
\ExplSyntaxOff

\zhvtset{
  grid color       = sealbrown,
  main font        = GenRyuMin2TW-SB,
  main font path   = C:/Users/Administrator/texmf/fonts/opentype/public/GenRyuMin2TW/,
jiazhu font      = GenRyuMin2TW-M,
  jiazhu font path = C:/Users/Administrator/texmf/fonts/opentype/public/GenRyuMin2TW/,
  % 页码框容纳 4 字：正常页「九十七」+前置空 1 格=4 行恰好；
  % 100+ 页「一百零一」共 5 行超出，末字「一/二/三」被框裁掉（2026-10-04 实测），故放宽到 5
  shukou page number chars = 5,
  % 书口书名框高从默认 6 字减到 4 字，把省出的 2 字高拨给 chapter 框，
  % 使 6 字书名+卷名放进 9 字高的 chapter 框后不再压线（2026-10-05）
  shukou title chars = 4,
}

\title{史记}[伍子胥列传]
\author{汉·司马迁}
\maker{闲叟漫记·重製}


% ===== 索隐述赞样式（2026-11-08 四次定稿 · 愚大人亲手改定参数）=====
% ⚠ 竖排前提：自定义字体家族必须带 RawFeature={vertical:+vert}；
%   simhei.ttf 加 +vhal 会整体加载失败静默回退，只写 +vert；simfang 同理。
\newCJKfontfamily{\shuzanfont}[Path=C:/Windows/Fonts/,RawFeature={vertical:+vert}]{simhei.ttf}
\newCJKfontfamily{\shuzanbody}[Path=C:/Windows/Fonts/,RawFeature={vertical:+vert}]{simfang.ttf}
% 题头：152x30pt 圆角 6pt 棕底白字；baseline (0,2) 使整框落入界行格子不压版框线；
% 白字节点 x 不补偿、y +8.2pt（像素级墨迹实测 dx/dy≈0，L/R 留白约 20pt）。
\newcommand{\shuzanhead}[1]{%
  \begin{tikzpicture}[baseline={(0,2)}]
    \fill[sealbrown, rounded corners=6pt]
      (-76pt,-15pt) rectangle (76pt,15pt);
    \node[text=white, inner sep=0pt, anchor=center] at (0pt,8.2pt)
      {\shuzanfont\fontsize{23}{23}\selectfont #1};
  \end{tikzpicture}%
}
% ★ 修正书口卷名位置：chapter 框起始位置改回 cls 默认值，避免全局影响。
%   只在每个卷首页通过 \juan 命令里的空命令微调该页位置（安全，不会跨页污染）。
\ExplSyntaxOn
\AtBeginDocument{
  \dim_set:Nn \l__zhvt_shukou_chapter_title_start_dim
    { \l__zhvt_shukou_title_height_dim + \l__zhvt_font_size_dim * 2 }
  \dim_set:Nn \l__zhvt_shukou_page_number_start_dim
    { \l__zhvt_shukou_chapter_title_start_dim + \l__zhvt_shukou_chapter_title_height_dim }
}
\ExplSyntaxOff

\begin{document}

% ★★ 版心（书口）排布覆写：把书名放到鱼尾之下（2026-10-04）
%
%   cls 默认版心自上而下：书名（title 静态框，在鱼尾之上）→ 鱼尾 → 篇名（chapter 动态框）→ 页码
%   本模板改为：          鱼尾（位置不动）→ 书名 → 篇名 → 页码分割线 → 页码
%
%   手法：把 cls 放在 title 静态框里的书名，改放进 chapter 动态框（鱼尾之下）；
%         title 静态框留空。书名取 \l__zhvt_title_str，即 \title{...} 的第一参数，
%         所以换书时只改 \title，不必在此重复写书名。
%
%   ⚠ 三条必须遵守的写法（2026-10-04 实测）：
%     ① 汉字必须用 \hbox{...} 包住。flowfram 会把框内容装进一个宽 14.4mm（= 一个 line_sep
%        = 书口栏宽）的纵向 minipage 并按该宽度折行，而整页靠 PDF /Rotate 90 显示、
%        框内文字本身不变向 —— 不包 \hbox 时「围炉夜话」4 字会被逐字折成 4 行，
%        在版面上就成了横排 4 字并横向溢出版框。
%     ② 竖排下**一个 \hbox 就是一整列**。书名与篇名要同处一列，必须放进同一个 \hbox；
%        分成两个 \hbox 时，后一个会被当成新的一列、甩到隔壁栏去
%        （实测「序」跑到书口栏外侧 14.4mm 的正文列上）。中间的 \quad 在框内是横向空白，
%        显示出来正好是书名与篇名之间的纵向间隔。
%     ③ ★ 内容盒要写成 **与框等宽的定宽盒** `\hbox to \l__zhvt_shukou_chapter_title_height_dim{…\hfil}`。
%        实测：chapter 框是把内容盒沿纵向**居中**摆放的，内容宽度一变，书名就跟着漂：
%        自然宽度的 \hbox（书名 38.4mm）落在 x=157.2mm，加篇名后（57.6mm）落在 x=152.4mm，
%        目录页（无篇名）与正文页的书名就差半格。做成等宽盒后，无论有没有篇名，
%        内容固定从框顶起排，实测两页均落在 x=142.8mm，不再漂。
\ExplSyntaxOn
% 把书口 chapter 框（书名+篇名）起始位置下移 1 字高，避免 4 字卷名顶到鱼尾
% ⚠ cls 的 \__zhvt_calculate_dimensions: 在 \AtBeginDocument 里执行，会重算该 dim；
%   必须把覆写挂到 begin document 之后，否则会被二次覆盖（2026-10-05 实测）。
% (旧方案已移到导言区，直接修改 chapter 框起始位置)
% 卷标题命令 #1=书口篇名(鱼尾) #2=正文卷名（含卷序，卷名内空格用 ~）
%
% ★★ 目录页两项修订（2026-10-05，愚大人指出）★★
%   ① 页码起点统一对齐：
%      原状 —— 目录项写成 \par\hspace{4em}卷名\hspace{4em}页码，页码起点 = 卷名实宽 + 4em。
%      卷名实宽随字数变：「卷一~梁惠王上」6 字 +~、「卷七~离娄上」5 字 +~，
%      差 1 个字位 → 卷七~卷十 4 条的页码整体比其余 10 条高 27.21pt（一个字），看着上下不齐。
%      改法 —— 卷名补全角空格（宽 = 1 字 = 27.215pt）补足到 **7 字符位**再写进目录，
%      14 条页码起点遂全部落在同一条竖线上（起点对齐；末字随页码字数自然参差）。
%      ⚠ 只对目录里的副本 \l_tmpc_tl 补，正文卷名仍用 \l_tmpa_tl 原样 —— 正文像素不受影响。
%   ② 卷名整体上移半字位：前导 \hspace 由 4em 改 3.5em。
%      实测原状目录块在版框内偏下（天头留白 139.6pt / 地脚留白 110.7pt），
%      上移 0.5em = 13.6pt 后上下基本均分，视觉居中。
\NewDocumentCommand { \juan } { mm } {
	\clearpage
	\tl_set:Nn \l__zhvt_chapter_title_tl { #1 }
	\tl_set:Nn \l_tmpa_tl { #2 }
	\tl_set_eq:NN \l_tmpc_tl \l_tmpa_tl
	\int_step_inline:nn { 7 - \str_count:n { #2 } }
	  { \tl_put_right:Nn \l_tmpc_tl { 　 } }
	\tl_set:Nx \l_tmpb_tl { \zhnum { page } }
	\tl_put_left:Nn \l_tmpb_tl { \hspace { 4em } }
	\tl_put_left:Nn \l_tmpb_tl { \l_tmpc_tl }
	\tl_put_left:Nn \l_tmpb_tl { \par \hspace { 3.5em } }
	\addtocontents { toc } { \l_tmpb_tl }
	\skip_horizontal:n { \l__zhvt_font_size_dim * 4 }
	\l_tmpa_tl
	\par
	\thispagestyle { empty }
}
\cs_set_protected:Npn \__zhvt_set_shukou_content: {
	% ★ 鱼尾上方留白（title 静态框）嵌入闲叟钤印：
	%   整页 /Rotate 90，故此处把图用 \rotatebox{90} 转正为竖章；
	%   尺寸卡在版心条幅内（书内约 7mm 宽 × 33.8mm 高，留余量不压鱼尾）。
	\setstaticcontents * { title } {
		\vbox_to_ht:nn { \l__zhvt_baseline_skip_dim } {
			\vskip 14.2pt
			\hbox_to_wd:nn { \l__zhvt_shukou_title_height_dim } {
				\hfil
				\rotatebox [ origin = c ] { 90 } { \includegraphics [ width = 7mm , height = 33.8mm ] { yshs.png } }
				\hfil
			}
		}
	}
% ★ 版心（书口）统一改法（2026-11-16 愚大人定）：18pt＋框内居中＋书口文字同版框色
%   ① 原 27.2pt 时 7 字长篇名（如「廉颇蔺相如列传」）末字溢出 chapter 框、压住页码；
%      改 \fontsize{18pt}{18pt}，前后 \hfil 令字串在框内居中（不随篇名字数漂移）。
%   ② \vbox{\vskip 8.2pt} 定量下移，抵消缩字号后栏宽方向 8.2pt 上偏（像素实测 y 中心 543.3 = 原版）。
%      ⚠ \raisebox / \vspace* 在 flowfram 框内无效；换字号须同步换补偿：16pt→10.1、19pt→7.4。
%   ③ chapter 与 page~number 均置 \color{sealbrown}（HTML 874434），与版框／界行同色；正文仍黑。
%   ⚠ 页码字位、钤印、鱼尾、界行、正文一律未动；五部合订本同此口径。
	\setdynamiccontents * { chapter }
	{ \color{sealbrown} \vbox { \vskip 8.2pt \hbox to \l__zhvt_shukou_chapter_title_height_dim
		{ \hfil \fontsize{18pt}{18pt}\selectfont \l__zhvt_title_str \quad \l__zhvt_chapter_title_tl \hfil } } }
	\setdynamiccontents * { page~number } { \color{sealbrown} \quad \zhnum { page } }
}
\__zhvt_set_shukou_content:
\ExplSyntaxOff

\maketitle

% 单篇，无多卷目录，不渲染目录页

\juan{伍子胥列传}{卷六十六~伍子胥列传第六}
\hspace*{2em}伍子胥者。楚人也。名员。员父曰伍奢。员兄曰伍尚。其先曰伍举。以直谏事楚庄王。有显。故其後世有名於楚。\\
\hspace*{2em}楚平王有太子名曰建。使伍奢为太傅。费无忌为少傅。无忌不忠於太子建。平王使无忌为太子取妇於秦。秦女好。无忌驰归报平王曰。﹁秦女绝美。王可自取。而更为太子取妇。﹂平王遂自取秦女而绝爱幸之。生子轸。更为太子取妇。\\
\hspace*{2em}无忌既以秦女自媚於平王。因去太子而事平王。恐一旦平王卒而太子立。杀己。乃因谗太子建。建母。蔡女也。无宠於平王。平王稍益疏建。使建守城父。备边兵。\\
\hspace*{2em}顷之。无忌又日夜言太子短於王曰。﹁太子以秦女之故。不能无怨望。原王少自备也。自太子居城父。将兵。外交诸侯。且欲入为乱矣。﹂平王乃召其太傅伍奢考问之。伍奢知无忌谗太子於平王。因曰。﹁王独柰何以谗贼小臣疏骨肉之亲乎。﹂无忌曰。﹁王今不制。其事成矣。王且见禽。﹂於是平王怒。囚伍奢。而使城父司马奋扬往杀太子。行未至。奋扬使人先告太子。﹁太子急去。不然将诛。﹂太子建亡奔宋。\\
\hspace*{2em}无忌言於平王曰。﹁伍奢有二子。皆贤。不诛且为楚忧。可以其父质而召之。不然且为楚患。﹂王使使谓伍奢曰。﹁能致汝二子则生。不能则死。﹂伍奢曰。﹁尚为人仁。呼必来。员为人刚戾忍卼。能成大事。彼见来之并禽。其势必不来。﹂王不听。使人召二子曰。﹁来。吾生汝父。不来。今杀奢也。﹂伍尚欲往。员曰。﹁楚之召我兄弟。非欲以生我父也。恐有脱者後生患。故以父为质。诈召二子。二子到。则父子俱死。何益父之死。往而令雠不得报耳。不如奔他国。借力以雪父之耻。俱灭。无为也。﹂伍尚曰。﹁我知往终不能全父命。然恨父召我以求生而不往。後不能雪耻。终为天下笑耳。﹂谓员。﹁可去矣。汝能报杀父之雠。我将归死。﹂尚既就执。使者捕伍胥。伍胥贯弓执矢乡使者。使者不敢进。伍胥遂亡。闻太子建之在宋。往从之。奢闻子胥之亡也。曰。﹁楚国君臣且苦兵矣。﹂伍尚至楚。楚并杀奢与尚也。\\
\hspace*{2em}伍胥既至宋。宋有华氏之乱。乃与太子建俱奔於郑。郑人甚善之。太子建又適晋。晋顷公曰。﹁太子既善郑。郑信太子。太子能为我内应。而我攻其外。灭郑必矣。灭郑而封太子。﹂太子乃还郑。事未会。会自私欲杀其从者。从者知其谋。乃告之於郑。郑定公与子产诛杀太子建。建有子名胜。伍胥惧。乃与胜俱奔吴。到昭关。昭关欲执之。伍胥遂与胜独身步走。几不得脱。追者在後。至江。江上有一渔父乘船。知伍胥之急。乃渡伍胥。伍胥既渡。解其剑曰。﹁此剑直百金。以与父。﹂父曰。﹁楚国之法。得伍胥者赐粟五万石。爵执珪。岂徒百金剑邪。﹂不受。伍胥未至吴而疾。止中道。乞食。至於吴。吴王僚方用事。公子光为将。伍胥乃因公子光以求见吴王。\\
\hspace*{2em}久之。楚平王以其边邑锺离与吴边邑卑梁氏俱蚕。两女子争桑相攻。乃大怒。至於两国举兵相伐。吴使公子光伐楚。拔其锺离。居巢而归。伍子胥说吴王僚曰。﹁楚可破也。原复遣公子光。﹂公子光谓吴王曰。﹁彼伍胥父兄为戮於楚。而劝王伐楚者。欲以自报其雠耳。伐楚未可破也。﹂伍胥知公子光有内志。欲杀王而自立。未可说以外事。乃进专诸於公子光。退而与太子建之子胜耕於野。\\
\hspace*{2em}五年而楚平王卒。初。平王所夺太子建秦女生子轸。及平王卒。轸竟立为後。是为昭王。吴王僚因楚丧。使二公子将兵往袭楚。楚发兵绝吴兵之後。不得归。吴国内空。而公子光乃令专诸袭刺吴王僚而自立。是为吴王阖庐。阖庐既立。得志。乃召伍员以为行人。而与谋国事。\\
\hspace*{2em}楚诛其大臣郤宛。伯州犁。伯州犁之孙伯嚭亡奔吴。吴亦以嚭为大夫。前王僚所遣二公子将兵伐楚者。道绝不得归。後闻阖庐弑王僚自立。遂以其兵降楚。楚封之於舒。阖庐立三年。乃兴师与伍胥。伯嚭伐楚。拔舒。遂禽故吴反二将军。因欲至郢。将军孙武曰。﹁民劳。未可。且待之。﹂乃归。\\
\hspace*{2em}四年。吴伐楚。取六与灊。五年。伐越。败之。六年。楚昭王使公子囊瓦将兵伐吴。吴使伍员迎击。大破楚军於豫章。取楚之居巢。\\
\hspace*{2em}九年。吴王阖庐谓子胥。孙武曰。﹁始子言郢未可入。今果何如。﹂二子对曰。﹁楚将囊瓦贪。而唐。蔡皆怨之。王必欲大伐之。必先得唐。蔡乃可。﹂阖庐听之。悉兴师与唐。蔡伐楚。与楚夹汉水而陈。吴王之弟夫概将兵请从。王不听。遂以其属五千人击楚将子常。子常败走。奔郑。於是吴乘胜而前。五战。遂至郢。己卯。楚昭王出奔。庚辰。吴王入郢。\\
\hspace*{2em}昭王出亡。入云梦。盗击王。王走郧。郧公弟怀曰。﹁平王杀我父。我杀其子。不亦可乎。﹂郧公恐其弟杀王。与王奔随。吴兵围随。谓随人曰。﹁周之子孙在汉川者。楚尽灭之。﹂随人欲杀王。王子綦匿王。己自为王以当之。随人卜与王於吴。不吉。乃谢吴不与王。\\
\hspace*{2em}始伍员与申包胥为交。员之亡也。谓包胥曰。﹁我必覆楚。﹂包胥曰。﹁我必存之。﹂及吴兵入郢。伍子胥求昭王。既不得。乃掘楚平王墓。出其尸。鞭之三百。然後已。申包胥亡於山中。使人谓子胥曰。﹁子之报雠。其以甚乎。吾闻之。人众者胜天。天定亦能破人。今子故平王之臣。亲北面而事之。今至於僇死人。此岂其无天道之极乎。﹂伍子胥曰。﹁为我谢申包胥曰。吾日暮途远。吾故倒行而逆施之。﹂於是申包胥走秦告急。求救於秦。秦不许。包胥立於秦廷。昼夜哭。七日七夜不绝其声。秦哀公怜之。曰。﹁楚虽无道。有臣若是。可无存乎。﹂乃遣车五百乘救楚击吴。六月。败吴兵於稷。会吴王久留楚求昭王。而阖庐弟夫概乃亡归。自立为王。阖庐闻之。乃释楚而归。击其弟夫概。夫概败走。遂奔楚。楚昭王见吴有内乱。乃复入郢。封夫概於堂谿。为堂谿氏。楚复与吴战。败吴。吴王乃归。\\
\hspace*{2em}後二岁。阖庐使太子夫差将兵伐楚。取番。楚惧吴复大来。乃去郢。徙於鄀。当是时。吴以伍子胥。孙武之谋。西破彊楚。北威齐晋。南服越人。\\
\hspace*{2em}其後四年。孔子相鲁。\\
\hspace*{2em}後五年。伐越。越王勾践迎击。败吴於姑苏。伤阖庐指。军卻。阖庐病创将死。谓太子夫差曰。﹁尔忘勾践杀尔父乎。﹂夫差对曰。﹁不敢忘。﹂是夕。阖庐死。夫差既立为王。以伯嚭为太宰。习战射。二年後伐越。败越於夫湫。越王勾践乃以馀兵五千人栖於会稽之上。使大夫种厚币遗吴太宰嚭以请和。求委国为臣妾。吴王将许之。伍子胥谏曰。﹁越王为人能辛苦。今王不灭。後必悔之。﹂吴王不听。用太宰嚭计。与越平。\\
\hspace*{2em}其後五年。而吴王闻齐景公死而大臣争宠。新君弱。乃兴师北伐齐。伍子胥谏曰。﹁勾践食不重味。吊死问疾。且欲有所用之也。此人不死。必为吴患。今吴之有越。犹人之有腹心疾也。而王不先越而乃务齐。不亦谬乎。﹂吴王不听。伐齐。大败齐师於艾陵。遂威邹鲁之君以归。益疏子胥之谋。\\
\hspace*{2em}其後四年。吴王将北伐齐。越王勾践用子贡之谋。乃率其众以助吴。而重宝以献遗太宰嚭。太宰嚭既数受越赂。其爱信越殊甚。日夜为言於吴王。吴王信用嚭之计。伍子胥谏曰。﹁夫越。腹心之病。今信其浮辞诈伪而贪齐。破齐。譬犹石田。无所用之。且盘庚之诰曰。﹃有颠越不恭。劓殄灭之。俾无遗育。无使易种于兹邑。﹄此商之所以兴。原王释齐而先越。若不然。後将悔之无及。﹂而吴王不听。使子胥於齐。子胥临行。谓其子曰。﹁吾数谏王。王不用。吾今见吴之亡矣。汝与吴俱亡。无益也。﹂乃属其子於齐鲍牧。而还报吴。\\
\hspace*{2em}吴太宰嚭既与子胥有隙。因谗曰。﹁子胥为人刚暴。少恩。猜贼。其怨望恐为深祸也。前日王欲伐齐。子胥以为不可。王卒伐之而有大功。子胥耻其计谋不用。乃反怨望。而今王又复伐齐。子胥专愎彊谏。沮毁用事。徒幸吴之败以自胜其计谋耳。今王自行。悉国中武力以伐齐。而子胥谏不用。因辍谢。详病不行。王不可不备。此起祸不难。且嚭使人微伺之。其使於齐也。乃属其子於齐之鲍氏。夫为人臣。内不得意。外倚诸侯。自以为先王之谋臣。今不见用。常鞅鞅怨望。原王早图之。﹂吴王曰。﹁微子之言。吾亦疑之。﹂乃使使赐伍子胥属镂之剑。曰。﹁子以此死。﹂伍子胥仰天叹曰。﹁嗟乎。谗臣嚭为乱矣。王乃反诛我。我令若父霸。自若未立时。诸公子争立。我以死争之於先王。几不得立。若既得立。欲分吴国予我。我顾不敢望也。然今若听谀臣言以杀长者。﹂乃告其舍人曰。﹁必树吾墓上以梓。令可以为器。而抉吾眼县吴东门之上。以观越寇之入灭吴也。﹂乃自刭死。吴王闻之大怒。乃取子胥尸盛以鸱夷革。浮之江中。吴人怜之。为立祠於江上。因命曰胥山。\\
\hspace*{2em}吴王既诛伍子胥。遂伐齐。齐鲍氏杀其君悼公而立阳生。吴王欲讨其贼。不胜而去。其後二年。吴王召鲁卫之君会之橐皋。其明年。因北大会诸侯於黄池。以令周室。越王勾践袭杀吴太子。破吴兵。吴王闻之。乃归。使使厚币与越平。後九年。越王勾践遂灭吴。杀王夫差。而诛太宰嚭。以不忠於其君。而外受重赂。与己比周也。\\
\hspace*{2em}伍子胥初所与俱亡故楚太子建之子胜者。在於吴。吴王夫差之时。楚惠王欲召胜归楚。叶公谏曰。﹁胜好勇而阴求死士。殆有私乎。﹂惠王不听。遂召胜。使居楚之边邑鄢。号为白公。白公归楚三年而吴诛子胥。\\
\hspace*{2em}白公胜既归楚。怨郑之杀其父。乃阴养死士求报郑。归楚五年。请伐郑。楚令尹子西许之。兵未发而晋伐郑。郑请救於楚。楚使子西往救。与盟而还。白公胜怒曰。﹁非郑之仇。乃子西也。﹂胜自砺剑。人问曰。﹁何以为。﹂胜曰。﹁欲以杀子西。﹂子西闻之。笑曰。﹁胜如卵耳。何能为也。﹂\\
\hspace*{2em}其後四岁。白公胜与石乞袭杀楚令尹子西。司马子綦於朝。石乞曰。﹁不杀王。不可。﹂乃劫王如高府。石乞从者屈固负楚惠王亡走昭夫人之宫。叶公闻白公为乱。率其国人攻白公。白公之徒败。亡走山中。自杀。而虏石乞。而问白公尸处。不言将亨。石乞曰。﹁事成为卿。不成而亨。固其职也。﹂终不肯告其尸处。遂亨石乞。而求惠王复立之。\\
\hspace*{2em}太史公曰。怨毒之於人甚矣哉。王者尚不能行之於臣下。况同列乎。向令伍子胥从奢俱死。何异蝼蚁。弃小义。雪大耻。名垂於後世。悲夫。方子胥窘於江上。道乞食。志岂尝须臾忘郢邪。故隐忍就功名。非烈丈夫孰能致此哉。白公如不自立为君者。其功谋亦不可胜道者哉。\\

% ===== 索隐述赞（司马贞《史记索隐》）· 统一样式 =====
% 题头「索隐述赞曰」五字并入反白框；正文每四句一列，定宽盒 \hfil 居列中；
% \fontsize{24}{41} 列距对齐全书界行（121+41k）；\clearpage 使述赞独占一页。
\clearpage
\shuzanhead{索隐述赞曰}\linebreak
{\shuzanbody\fontsize{24}{41}\selectfont
\hbox to \hsize{\hfil 谗人罔极。交乱四国。嗟彼伍氏。被兹凶慝。\hfil}\linebreak
\hbox to \hsize{\hfil 员独忍诟。志复冤毒。霸吴起师。伐楚逐北。\hfil}\linebreak
\hbox to \hsize{\hfil 鞭尸雪耻。抉眼弃德。\hfil}%
\par}
\end{document}
